\documentclass[10pt,a4paper]{amsart}
\usepackage[margin=19mm]{geometry}
\usepackage{amsmath,amssymb,mathtools}
\usepackage[scr=rsfs,cal=euler]{mathalfa}
\usepackage{xfrac}
\usepackage[unicode,hidelinks,bookmarks=false]{hyperref}
\newcommand{\ie}{i.e.\ }
\newcommand{\cf}{cf.\ }
\newcommand{\resp}{respectively\ }
\newcommand{\Subsection}{\subsection}
\providecommand{\nobreakdash}{\nobreak\hbox{-}}
\setlength{\emergencystretch}{2em}
\multlinegap0pt
\usepackage{amssymb}
\usepackage{tikz}
\newtheorem{lemma}{Lemma}
\title{Detecting Zariski Pairs by Algorithms and Computational Classification in Conic--Line Arrangements}
\author{Meirav Amram, Gal Goren}
\date{Source: arXiv:2601.00463v1 (CC BY 4.0)}
\begin{document}
\maketitle

\noindent\emph{Source and licence.} Detecting Zariski Pairs by Algorithms and Computational Classification in Conic--Line Arrangements. Version arXiv:2601.00463v1 (CC BY 4.0), distributed by arXiv under the Creative Commons Attribution 4.0 International licence, http://creativecommons.org/licenses/by/4.0/, as recorded in the arXiv licence metadata for this exact version and shown on the abstract page https://arxiv.org/abs/2601.00463v1. Full licence notice: \texttt{licenses/arxiv-2601.00463-CC-BY-4.0.txt}.\par\smallskip

\noindent\textit{Editorial scope.} Counted unit: the article's lemma lemma2 (source0.tex lines 423--426). The article's complete original proof of it is delivered with it (source0.tex lines 428--441, one \texttt{proof} environment of 14 lines). No mathematics is written, repaired or re-derived: the statement and the proof are the article's own lines, in the article's own notation, and no retained line is substituted or re-flowed.  No further source-local context is needed: every cross-reference of every retained line resolves inside the two delivered spans.  Nothing was omitted from any retained span, so the package carries no omission record.  No retained line contains a citation, so the wrapper carries no bibliography and no bibliography entry.  No retained line contains a figure environment, an image-inclusion command or drawing code, so no image is delivered and the package declares no asset dependency.  Editorial formatting only, in addition: an AMS \texttt{amsart} preamble that restates the article's own declarations and macros used by the retained lines, namely the environments \texttt{lemma} on separate counters, together with the standard packages those lines need.  The retained environments print in this document as Lemma 1 in order of appearance, whereas the article prints the same items with the numbering of its own full document.  The complete native source of the article as distributed in the arXiv e-print for this exact version is bundled unchanged as \texttt{sources/192001/source0.tex}.
\begin{lemma}\label{lemma2}
If a non-tangent line $L$ has exactly one constraint $p$, then the characteristic of $p$ in 
$\overline{B_1 \setminus L}$ occurs at least at one other point in $\overline{B_1 \setminus L}$.
\end{lemma}

\begin{proof}
Assume for contradiction that the characteristic of $p$ is unique in 
$\overline{B_1 \setminus L}$.  
Let $L'$ be the line in $B_2$ corresponding to $L$, so that 
$\overline{B_1 \setminus L} \simeq \overline{B_2 \setminus L'}$.  
Since the characteristic of $p$ in $\overline{B_1 \setminus L}$ is unique, and $L$ is not tangent to the conic, there is a unique way, up to a continuous deformation that preserves the combinatorics, to add a line to $\overline{B_1 \setminus L}$ to recover 
the combinatorics of $B_1$.  
Thus
\[
B_1 = \overline{B_1 \setminus L} \cup L \;\;\simeq\;\; 
\overline{B_2 \setminus L'} \cup L' = B_2,
\]
which is a contradiction.
\end{proof}

\end{document}
